\documentclass{article}
\usepackage{amsmath}

\title{Rotation Matrices with Euler Angles}
\author{}
\date{}

\begin{document}

\maketitle

\section{Introduction}
Euler angles define the orientation of a rigid body using three consecutive rotations about different coordinate axes. A common convention is the ZYX sequence, involving:
\begin{itemize}
    \item Yaw ($R_z(\alpha)$): Rotation around the $z$-axis.
    \item Pitch ($R_y(\beta)$): Rotation around the $y$-axis.
    \item Roll ($R_x(\gamma)$): Rotation around the $x$-axis.
\end{itemize}

\section{Rotation Matrices}
Each individual rotation matrix is given by:

\subsection{Rotation About the $x$-Axis (Roll)}
\begin{equation}
R_x(\gamma) =
\begin{bmatrix}
1 & 0 & 0 \\
0 & \cos\gamma & -\sin\gamma \\
0 & \sin\gamma & \cos\gamma
\end{bmatrix}
\end{equation}

\subsection{Rotation About the $y$-Axis (Pitch)}
\begin{equation}
R_y(\beta) =
\begin{bmatrix}
\cos\beta & 0 & \sin\beta \\
0 & 1 & 0 \\
-\sin\beta & 0 & \cos\beta
\end{bmatrix}
\end{equation}

\subsection{Rotation About the $z$-Axis (Yaw)}
\begin{equation}
R_z(\alpha) =
\begin{bmatrix}
\cos\alpha & -\sin\alpha & 0 \\
\sin\alpha & \cos\alpha & 0 \\
0 & 0 & 1
\end{bmatrix}
\end{equation}

\section{Combined Rotation Matrix (ZYX Sequence)}
The total rotation matrix is given by:
\begin{equation}
R(\alpha, \beta, \gamma) = R_z(\alpha) R_y(\beta) R_x(\gamma)
\end{equation}
Expanding this product:
\begin{equation}
R(\alpha, \beta, \gamma) =
\begin{bmatrix}
\cos\alpha \cos\beta & \cos\alpha \sin\beta \sin\gamma - \sin\alpha \cos\gamma & \cos\alpha \sin\beta \cos\gamma + \sin\alpha \sin\gamma \\
\sin\alpha \cos\beta & \sin\alpha \sin\beta \sin\gamma + \cos\alpha \cos\gamma & \sin\alpha \sin\beta \cos\gamma - \cos\alpha \sin\gamma \\
-\sin\beta & \cos\beta \sin\gamma & \cos\beta \cos\gamma
\end{bmatrix}
\end{equation}
\section{Rotational Velocity Matrix}
The angular velocity vector $\boldsymbol{\omega}$ in terms of Euler angles is given by:
\begin{equation}
\boldsymbol{\omega} = 
\begin{bmatrix}
\dot{\alpha} \\
0 \\
0
\end{bmatrix} + R_x(\gamma)^T 
\begin{bmatrix}
0 \\
\dot{\beta} \\
0
\end{bmatrix} + R_x(\gamma)^T R_y(\beta)^T 
\begin{bmatrix}
0 \\
0 \\
\dot{\gamma}
\end{bmatrix}
\end{equation}
Expanding this, the rotational velocity matrix $\Omega$ is:
\begin{equation}
\Omega = 
\begin{bmatrix}
1 & 0 & -\sin\beta \\
0 & \cos\gamma & \cos\beta \sin\gamma \\
0 & -\sin\gamma & \cos\beta \cos\gamma
\end{bmatrix}
\begin{bmatrix}
\dot{\alpha} \\
\dot{\beta} \\
\dot{\gamma}
\end{bmatrix}
\end{equation}
\mathbf{\Omega} =
\begin{pmatrix}
\omega_1 \\
\omega_2 \\
\omega_3
\end{pmatrix}
=
\begin{pmatrix}
\dot{\phi} \sin\theta \sin\psi + \dot{\theta} \cos\psi \\
\dot{\phi} \sin\theta \cos\psi + \dot{\theta} \sin\psi \\
\dot{\psi} + \dot{\phi} \cos\theta
\end{pmatrix}

\text{סוגרי פואסון} 

\{f,g\}_{pq} := \sum \left( \frac{\partial f}{\partial p_i} \frac{\partial g}{\partial q_i} - \frac{\partial f}{\partial q_i} \frac{\partial g}{\partial p_i} \right)

\{f,g\} = -\{g,f\} , \quad \{a f + b g, h\} = a \{f,h\} + b \{g,h\}

\{q_i, q_j\} = \{p_i, p_j\} = 0, \quad \{p_i, q_j\} = \delta_{ij}

\{a b, c\} = a \{b,c\} + b \{a,c\}

\{f, \{g,h\}\} + \{g, \{h,f\}\} + \{h, \{f,g\}\} = 0

\{f(u), g(v)\} = \{u,v\} \frac{\partial f}{\partial u} \frac{\partial g}{\partial v}

\{f(\vec{p}), g(\vec{q})\} = \sum \frac{\partial f}{\partial p_i} \frac{\partial g}{\partial q_i}

\{f(\vec{q}, \vec{p}), q_i\} = \frac{\partial f}{\partial p_i}, \quad \{ f(\vec{q}, \vec{p}), p_i\} = -\frac{\partial f}{\partial q_i}

\{L_i, L_j\} = -\sum_k \epsilon_{ijk} L_k

\section{Applications}
Euler angles are widely used in physics and engineering, including:
\begin{itemize}
    \item Aerospace: Aircraft orientation (pitch, yaw, roll).
    \item Robotics: Positioning robotic arms.
    \item Computer Graphics: 3D object rotation.
    \item Quantum Mechanics: Spin transformations.
\end{itemize}

\end{document}